\documentclass[10pt,a4paper]{amsart}
\usepackage[margin=19mm]{geometry}
\usepackage{amsmath,amssymb,mathtools}
\usepackage[scr=rsfs,cal=euler]{mathalfa}
\usepackage{xfrac}
\usepackage[unicode,hidelinks,bookmarks=false]{hyperref}
\newcommand{\ie}{i.e.\ }
\newcommand{\cf}{cf.\ }
\newcommand{\resp}{respectively\ }
\newcommand{\Subsection}{\subsection}
\providecommand{\nobreakdash}{\nobreak\hbox{-}}
\setlength{\emergencystretch}{2em}
\multlinegap0pt
\usepackage{amssymb}
\newtheorem{proposition}{Proposition}
\newtheorem{theorem}{Theorem}
\newtheorem{corollary}{Corollary}
\newtheorem{spacing}{Spacing}
\usepackage{cleveref}
\crefname{proposition}{Proposition}{Propositions}
\crefname{theorem}{Theorem}{Theorems}
\crefname{corollary}{Corollary}{Corollarys}
\crefname{spacing}{Spacing}{Spacings}
\newcommand{\GL}{\mathrm{GL}}
\newcommand{\bC}{\mathbb{C}}
\newcommand{\bP}{\mathbb{P}}
\newcommand{\calC}{\mathcal{C}}
\newcommand{\calO}{\mathcal{O}}
\title{Cubic fourfolds containing highly singular hyperplane sections}
\author{Lisa Marquand, Sasha Viktorova}
\date{Source: arXiv:2601.20778v1 (CC BY 4.0)}
\begin{document}
\maketitle

\noindent\emph{Source and licence.} Cubic fourfolds containing highly singular hyperplane sections. Version arXiv:2601.20778v1 (CC BY 4.0), distributed by arXiv under the Creative Commons Attribution 4.0 International licence, http://creativecommons.org/licenses/by/4.0/, as recorded in the arXiv licence metadata for this exact version and shown on the abstract page https://arxiv.org/abs/2601.20778v1. Full licence notice: \texttt{licenses/arxiv-2601.20778-CC-BY-4.0.txt}.\par\smallskip

\noindent\textit{Editorial scope.} Counted unit: the article's theorem divisors (source0.tex lines 343--347). The article's complete original proof of it is delivered with it (source0.tex lines 472--652, one \texttt{proof} environment of 181 lines, which the article places away from the statement). No mathematics is written, repaired or re-derived: the statement and the proof are the article's own lines, in the article's own notation, and no retained line is substituted or re-flowed.  Retained uncounted as the source-local context the proof needs: lines 310--325; lines 335--337; lines 421--423.  Nothing was omitted from any retained span, so the package carries no omission record.  No retained line contains a citation, so the wrapper carries no bibliography and no bibliography entry.  No retained line contains a figure environment, an image-inclusion command or drawing code, so no image is delivered and the package declares no asset dependency.  Editorial formatting only, in addition: an AMS \texttt{amsart} preamble that restates the article's own declarations and macros used by the retained lines, namely the environments \texttt{proposition}, \texttt{theorem}, \texttt{corollary}, \texttt{spacing} on separate counters, together with the standard packages those lines need.  The retained environments print in this document as Proposition 1, Theorem 1, Corollary 1, Spacing 1 in order of appearance, whereas the article prints the same items with the numbering of its own full document.  The complete native source of the article as distributed in the arXiv e-print for this exact version is bundled unchanged as \texttt{sources/193548/source0.tex}.
\begin{proposition} \label{prop: equations}
    Let $Y\subset \bP^4$ be a $K$-cubic threefold with $K\in \{E_6, D_6, D_5+A_1, D_4+A_2, T_{333}\}.$ After a linear changes of coordinates, we have that $Y=V(f_K)$, where $f_K$ is given as:
    \begin{align*}
        f_{E_6}&:=x_0x_1x_2+x_1q_1(x_1,x_3,x_4)+x_2q_2(x_2,x_3,x_4)+x_1x_2l(x_1,\ldots,x_4)+x_3^3,\\
        f_{D_6}&:=x_0x_1x_2+x_1q_1(x_1,x_3,x_4)+x_2q_2(x_2,x_3)+x_2x_4h(x_2,x_3)+x_1x_2l(x_1,\ldots,x_4)+x_3^2x_4,\\
        f_{D_5+A_1}&:=x_0x_1x_2+x_1q_1(x_1,x_3,x_4)+x_2q_2(x_3,x_4)+x_1x_2l(x_1,\ldots,x_4)+x_3^2x_4,\\
        f_{D_4+A_2}&:=x_0x_1x_2+x_1q_1(x_1,x_3,x_4)+x_2h^2(x_3,x_4)+x_1x_2l(x_1,\ldots,x_4)+x_3^2x_4+x_4^3,\\
        f_{T_{333}}&:=x_0x_1^2+f_3(x_1,\ldots,x_4).
    \end{align*}

    A $(D_4+2A_1)$-cubic threefold $Y\subset\bP^4$ with $\sigma(Y)=0$ or $\sigma(Y)=1$ can similarly be defined by:
    \begin{align*}
    f_{D_4+2A_1}^0&:=x_0x_1x_2+x_1q_1(x_3,x_4)+x_2q_2(x_3,x_4)+x_1x_2l(x_1,\ldots,x_4)+x_3^2x_4+x_4^3,\;\text{or}\\        f_{D_4+2A_1}^1&:=x_0x_1x_2+x_1q_1(x_1,x_3,x_4)+x_2x_4h(x_3,x_4)+x_1x_2l(x_1,\ldots,x_4)+x_3^2x_4+x_4^3.
    \end{align*}
    Here the polynomials $h$ and $l$ are linear, $q_1$ and $q_2$ are homogeneous of degree $2$, $f_3$ is homogeneous of degree $3$.
\end{proposition}

\begin{equation}\label{eqn: fourfold}
        F_K:=x_5q(x_0,\ldots,x_5)+f_K.
    \end{equation}

\begin{corollary} \label{E6 discrete fibers}
    A general smooth cubic fourfold containing a $K$-hyperplane section for $K\in \{E_6, D_6,D_5+A_1, D_4+A_2,D_4+2A_1\}$ has finitely many $K$-hyperplane sections.
\end{corollary}

\begin{theorem}\label{theorem: divisors}
    The loci $D_{K}\subset\calC$ are irreducible for $K\in\{E_6,D_6,D_5+A_1,D_4+A_2,T_{333}\}$. The locus $D_{D_4+2A_1}$ has two components, denoted $D_{D_4+2A_1}^0$ and $D_{D_4+2A_1}^1$, distinguished by whether the defect $\sigma$ is $0$ or $1$.
    
    The loci $D_{E_6}, D_{D_6}, D_{D_5+A_1}, D_{D_4+A_2}, D_{D_4+2A_1}^0, D_{D_4+2A_1}^1$ are irreducible divisors in $\calC$, and $D_{T_{333}}$ has codimension $2$ in $\calC$.
\end{theorem}

\begin{proof}[Proof of \Cref{theorem: divisors}]
    Let  $K\in\{E_6, D_6, D_5+A_1, D_4+A_2, T_{333}\}$ and consider the subset $W_K\subset \bP(H^0(\bP^5, \calO(3)))$ parametrising cubic fourfolds with equation of the form:
    \begin{equation}
        F_K:=x_5q(x_0,\ldots,x_5)+f_K=0
    \end{equation}
    as before.
    Notice that $W_K$ can be identified with its preimage in $H^0(\bP^5, \calO(3))$, as the Equation \ref{eqn: fourfold} contains an $x_0x_1x_2$ or $x_0x_1^2$ term with fixed coefficient $1$ for each $K$. 
    We abuse notation and call this subset also $W_K$. 
    By definition, the locus $D_{K}$ is the closure of the image in $\calC$ of an open subset of $W_K$. 
    Since $W_K$ is isomorphic to an affine space and thus is irreducible, $D_{K}$ is irreducible as well.

    \begin{spacing}{1.25}
    Similarly, let $W_{D_4+2A_1}^0,W_{D_4+2A_1}^1\subset H^0(\bP^5, \calO(3))$ be the loci parametrising cubic fourfolds given by$F_K=0$ with $f_K:=f^0_{D_4+2A_1},f^1_{D_4+2A_1}$ respectively, as in Proposition \ref{prop: equations}. 
    By the same argument each $W_{D_4+2A_1}^i$ can be identified with its image in $\bP(H^0(\bP^5,\calO(3)))$. 
    Each $W_{D_4+2A_1}^i$ is isomorphic to an affine space and is irreducible, thus each component $D_{D_4+2A_1}^0, D_{D_4+2A_1}^1$ is irreducible as well.
  \end{spacing}
  
    The quadratic form $q(x_0,\dots x_5)$ has $21$ parameters. The equation $f_K$ depends on $p_K$ parameters, where $p_K$ is as in the table below.

\begin{table}[h!]
\centering
\begin{tabular}{c|c|c|c|c|c|c|c}

$K$ & $E_6$ & $D_6$ & $D_5+A_1$& $D_4+A_2$ & $D_4+2A_1$, $\sigma=0$ & $D_4+2A_1$, $\sigma=1$ & $T_{333}$ \\ \hline
$p_K$   & 16   & 15   & 13   & 12  & 10  & 12 & 20 \\ 
\end{tabular}

\label{tab: parameter counts}
\end{table}
        
    Let $d_K$ be the dimension of the intersection of $W_K$ and the $\GL(6)$-orbit of a general element $w_K\in W_K$. 
    We have 
    \begin{equation}\label{eqn: dimension}
        \dim D_K=\dim W_K-\dim(W_K\cap (w\cdot\GL(6)))=(21+p_K)-d_K.
    \end{equation}
    We claim that $d_K$ is as in the following table:

    \begin{table}[h!]
\centering
\begin{tabular}{c|c|c|c|c|c|c|c}

$K$ & $E_6$ & $D_6$ & $D_5+A_1$& $D_4+A_2$ & $D_4+2A_1$, $\sigma=0$ & $D_4+2A_1$, $\sigma=1$ & $T_{333}$ \\ \hline
$d_K$   & 18   & 17   &  15  & 14   & 12   & 14 & 23\\ 
\end{tabular}

\label{tab: dK counts}
\end{table}
    
   Assuming the claim, the Theorem follows from Equation \ref{eqn: dimension}; it remains to compute $d_K$.
    
    \medskip
    \begin{spacing}{1.25}
    Consider the subgroup $G<\GL(6)$ consisting of elements preserving the plane $x_5=0$. By \Cref{E6 discrete fibers}, we have $\dim(W_K\cap w_K\cdot G)=\dim(W_K\cap w_K\cdot\GL(6))=d_K$.
    Indeed, let $Y_1,\ldots,Y_m$ be all the $K$-hyperplane sections of $X=V(w_K)\subset \bP^5$. There are suitable coordinate changes for each $i=1,\ldots,m$ such that $X$ is isomorphic to $X_i:=V(w_K^i)$ for $w_K^i\in W_K$ and $Y_i$ is given by $x_5=0$. Then  $W_K\cap w_K\cdot GL(6)=\cup_{i=1}^m(W_K\cap w_K^i\cdot G)$ and $\dim(W_K\cap w_K\cdot GL(6))=\dim(W_K\cap w_K^i\cdot G)$ for all $i=1,\ldots,m$ by the generality assumption.
    Notice that a general element $w_K\in W_K$ is smooth and thus GIT-semistable. Hence it has finite stabilizer and $d_K$ is equal to the dimension of $H_{w_K}:=\{g\in G \mid g\cdot w_K\in W_K\}$.
    \end{spacing}
    
    The group $G$ is isomorphic to the semidirect product $\GL(5)\ltimes(\bC^5\times \bC^*)$, where $\GL(5)$ is the subgroup of $\GL(6)$ acting on the first five coordinate functions $x_0,\ldots,x_4$. 
    The group $\bC^5\times \bC^*$ acts as the group of transformations of the form $x_i\mapsto \alpha_ix_5+x_i$, $x_5\mapsto \lambda x_5$, where $(\alpha_0,\ldots,\alpha_4,\lambda)\in \bC^5\times \bC^*$.
    
    We will now fix a general element $u_K$ in the set $U_K\subset H^0(\bP^4,\calO(3))$ of polynomials of the form $f_K\in \bC[x_0,\dots x_4]$ and describe the subset $H_{u_K}$ of $\GL(5)$ consisting of elements $h\in\GL(5)$ such that $u_K\cdot h$ is contained in $U_K$.  We will describe $H_{u_K}$ for each $K$ separately.  
    Notice that the unique singular point $p:=[1:0:0:0:0]$ of the threefold defined by $u_K$ has to be fixed. If $p$ is fixed, then an element in $H_{u_K}$ has the form: 
    $$x_0\mapsto a_{00}x_0+\ldots a_{04}x_4, \hspace{1cm} x_i\mapsto a_{i1}x_1+\ldots a_{i4}x_4, \text{ for } i=1,\dots 4.$$ For each $K$, we have $d_K=\dim(W_K\cap w_K\cdot G)=\dim(U_K\cap u_K\cdot\GL(5))+\dim(\bC^5\times \bC^*)=\dim H_{u_K}+6$.
    
      \smallskip
      
    \noindent\textbf{Case 1: }Let $K=E_6$. Recall that $u_{E_6}$ has the form:
    \[
        x_0x_1x_2+x_1q_1(x_1,x_3,x_4)+ x_2q_2(x_2,x_3,x_4)+x_1x_2l(x_1,\ldots,x_4)+x_3^3.
    \]
    As there is exactly one term in $u_{E_6}$ containing $x_0$, the subset $H_{u_K}\subset \GL(5)$ can only contain elements mapping $x_1\mapsto a_{11}x_1$ and $x_2\mapsto a_{22}x_2$, or $x_1\mapsto a_{12}x_2$ and $x_2\mapsto a_{21}x_1$, i.e., the variables $x_1$ and $x_2$ can only be rescaled and switched. 
    Moreover, since the term $x_0x_1x_2$ has a fixed coefficient $1$, it follows that $a_{00}a_{11}a_{22}=1$ or $a_{00}a_{12}a_{21}=1$. 
    Also, we note that $a_{34}=0$ (since we cannot have a nontrivial coefficient if front of $x_4^3$) and $a_{33}$ is a cube root of unity. 
    It is easy to see that for a transformation $h$ satisfying these properties $u_{E_6}\cdot h\in U_{E_6}$, thus we get that $H_{u_{E_6}}$ is generated by matrices    
    $$
    \begin{pmatrix}
        a_{00} & a_{01} & a_{02} & a_{03} & a_{04}\\
             0 & a_{11} &      0 &      0 &      0\\
             0 &      0 & a_{22} &      0 &      0\\
             0 & a_{31} & a_{32} & a_{33} &      0\\
             0 & a_{41} & a_{42} & a_{43} & a_{44}
    \end{pmatrix} \text{ and }
    \begin{pmatrix}
             1 &      0 &      0 &      0 &      0\\
             0 &      0 &      1 &      0 &      0\\
             0 &      1 &      0 &      0 &      0\\
             0 &      0 &      0 &      1 &      0\\
             0 &      0 &      0 &      0 &      1
    \end{pmatrix}, \text{ where } a_{00}a_{11}a_{22}=1 \text{ and } a_{33}^3=1.
    $$
    The set $H_{u_{E_6}}$ is of dimension $12$. We see $d_{E_6}=\dim H_{u_{E_6}}+6=12+6=18$.

    \medskip    

    \noindent \textbf{Case 2:} In the $K=D_6$ case, we proceed similarly. Again, there is exactly one term in $u_{D_6}$ containing $x_0$, and we can only have $x_1\mapsto a_{11}x_1$ and $x_2\mapsto a_{22}x_2$, or $x_1\mapsto a_{12}x_2$ and $x_2\mapsto a_{21}x_1$; we also have $a_{00}a_{11}a_{22}=1$ or $a_{00}a_{12}a_{21}=1$. Since $x_3^2x_4$ is the only monomial of $u_{D_6}$ in variables $x_3$ and $x_4$, we conclude that $a_{34}=a_{43}=0$ and $a_{33}^2a_{44}=1$. Finally, we notice that the form of $u_{D_6}$ is not symmetric in $x_1$ and $x_2$, thus transformations with $x_1\mapsto a_{12}x_2$ and $x_2\mapsto a_{21}x_1$ cannot be in $H_{u_{D_6}}$. We see that $H_{u_{D_6}}$ consists of matrices
    $$
    \begin{pmatrix}
        a_{00} & a_{01} & a_{02} & a_{03} & a_{04}\\
             0 & a_{11} &      0 &      0 &      0\\
             0 &      0 & a_{22} &      0 &      0\\
             0 & a_{31} & a_{32} & a_{33} &      0\\
             0 & a_{41} & a_{42} &      0 & a_{44}
    \end{pmatrix}, \text{ where } a_{00}a_{11}a_{22}=1 \text{ and } a_{33}^2a_{44}=1.
    $$
    The set $H_{D_6}$ is of dimension $11$ and $d_{D_6}=11+6=17$.

    \medskip

    \noindent\textbf{Case 3:} In the $K=D_5+A_1$ case, we see that $H_{u_{D_5+A_1}}$ consists of matrices:
$$
\begin{pmatrix}
        a_{00} & a_{01} & a_{02} & a_{03} & a_{04}\\
        0 & a_{11} &      0 &      0 &      0\\
        0 &      0 & a_{22} &      0 &      0\\
        0 & a_{31} &      0 & a_{33} &      0\\
        0 & a_{41} &      0 &      0 & a_{44}
\end{pmatrix}, 
$$
where $a_{00}a_{11}a_{22}=1$ and $a_{33}^2a_{44}=1.$
The set of matrices is $9$-dimensional and $d_K=15$. 

\medskip

\noindent\textbf{Case 4:} In the case $K=D_4+A_2$, we see that $H_{u_{D_4+A_2}}$ consists of matrices:
$$
\begin{pmatrix}
        a_{00} & a_{01} & a_{02} & a_{03} & a_{04}\\
        0 & a_{11} &      0 &      0 &      0\\
        0 &      0 & a_{22} &      0 &      0\\
        0 & a_{31} &      0 & a_{33} & a_{34}\\
        0 & a_{41} &      0 & a_{43} & a_{44}
\end{pmatrix},$$
where $a_{00}a_{11}a_{22}=1$ and there are 18 solutions for $a_{33},a_{34},a_{43},a_{44}$. The set of matrices is $8$-dimensional and $d_K=14$. 

\medskip


\noindent\textbf{Case 5:} In the case $K=D_4+2A_1$ and $\sigma=0$, we see that $H_{u_{D_4+2A_1}}^0$ consists of matrices:
$$
\begin{pmatrix}
        a_{00} & a_{01} & a_{02} & a_{03} & a_{04}\\
        0 & a_{11} &      0 &      0 &      0\\
        0 &      0 & a_{22} &      0 &      0\\
        0 &      0 &      0 & a_{33} & a_{34}\\
        0 &      0 &      0 & a_{43} & a_{44}
\end{pmatrix},
\begin{pmatrix}
        1 &      0 &      0 &      0 &      0\\
        0 &      0 &      1 &      0 &      0\\
        0 &      1 &      0 &      0 &      0\\
        0 &      0 &      0 &      1 &      0\\
        0 &      0 &      0 &      0 &      1
\end{pmatrix},
$$
where $a_{00}a_{11}a_{22}=1$ and there are 18 solutions for $a_{33},a_{34},a_{43},a_{44}$.
The set of matrices is $6$-dimensional and $d_K=12$.
\medskip

\noindent \textbf{Case: 6} In the case $K=D_4+2A_1$ and $\sigma=1$, we see that $H_{u_{D_4+2A_1}}^1$ consists of matrices:
$$
\begin{pmatrix}
        a_{00} & a_{01} & a_{02} & a_{03} & a_{04}\\
        0 & a_{11} &      0 &      0 &      0\\
       0 &      0 & a_{22} &      0 &      0\\
       0 & a_{31} &      0 & a_{33} & a_{34}\\
      0 & a_{41} &      0 & a_{43} & a_{44}
\end{pmatrix}, $$
where $a_{00}a_{11}a_{22}=1$ and there are 18 solutions for $a_{33},a_{34},a_{43},a_{44}$. The set of matrices is $8$-dimensional and $d_K=14$. 
    
    \noindent \textbf{Case 7:} In the $K= T_{333}$ case, we can only have $x_1\mapsto \pm x_1$ and we see that $H_{u_{T_{333}}}$ consists of matrices
    $$
    \begin{pmatrix}
        a_{00} & a_{01} & a_{02} & a_{03} & a_{04}\\
             0 &  \pm 1 &      0 &      0 &      0\\
             0 & a_{21} & a_{22} & a_{23} & a_{24}\\
             0 & a_{31} & a_{32} & a_{33} & a_{34}\\
             0 & a_{41} & a_{42} & a_{43} & a_{44}
    \end{pmatrix},
    $$
    thus $d_{T_{333}}=17+6=23$, proving the claim.
\end{proof}

\end{document}
